\documentclass{article}
\usepackage{apacite}
\usepackage{hyperref}
\usepackage{multicol}
\usepackage{natbib}
\usepackage{sophia}
\usepackage[utf8]{inputenc}
\usepackage{array}
\setlength{\parindent}{0pt}
\usepackage{hyperref}
\usepackage{physics}
\usepackage{soul}
\makeatletter
\newcommand*{\rom}[1]{\expandafter\@slowromancap\romannumeral #1@}

\usetikzlibrary{calc,patterns,angles,quotes}
\usepackage{pgfplots}
\pgfplotsset{width=15cm,height=10cm,compat=1.9}

\pgfplotstableread{
x	y	delta-x	delta-y
-0.009	-6.3255	0.0284	0.0028
0.0611	-5.9201	0.0302	0.0019
0.2967	-5.6324	0.0375	0.0014
0.4689	-5.4092	0.0445	0.0011
0.4903	-5.2269	0.0455	0.0009
0.4822	-5.0728	0.0451	0.0008
0.5518	-4.9392	0.0485	0.0007
}{\mytable}

\begin{filecontents}{graphdata.dat}
x	y	delta-x	delta-y
-0.009	-6.3255	0.0284	0.0028
0.0611	-5.9201	0.0302	0.0019
0.2967	-5.6324	0.0375	0.0014
0.4689	-5.4092	0.0445	0.0011
0.4903	-5.2269	0.0455	0.0009
0.4822	-5.0728	0.0451	0.0008
0.5518	-4.9392	0.0485	0.0007
\end{filecontents}

\usepackage[english]{babel}
\usepackage{fancyhdr}
\usepackage{lastpage}
\pagestyle{fancy}
\fancyhf{}


\pagestyle{empty}% for cropping
\usepackage{tikz}
\usetikzlibrary{calc}
\usepackage{zref-savepos}
\usepackage{tabu}

\newcounter{NoTableEntry}
\renewcommand*{\theNoTableEntry}{NTE-\the\value{NoTableEntry}}

\newcommand*{\strike}[2]{%
  \multicolumn{1}{#1}{%
    \stepcounter{NoTableEntry}%
    \vadjust pre{\zsavepos{\theNoTableEntry t}}% top
    \vadjust{\zsavepos{\theNoTableEntry b}}% bottom
    \zsavepos{\theNoTableEntry l}% left
    \hspace{0pt plus 1filll}%
    #2% content
    \hspace{0pt plus 1filll}%
    \zsavepos{\theNoTableEntry r}% right
    \tikz[overlay]{%
      \draw
        let
          \n{llx}={\zposx{\theNoTableEntry l}sp-\zposx{\theNoTableEntry r}sp-\tabcolsep},
          \n{urx}={\tabcolsep},
          \n{lly}={\zposy{\theNoTableEntry b}sp-\zposy{\theNoTableEntry r}sp},
          \n{ury}={\zposy{\theNoTableEntry t}sp-\zposy{\theNoTableEntry r}sp}
        in
        (\n{llx}, \n{lly}) -- (\n{urx}, \n{ury})
      ;
    }% 
  }%
}

\usepackage{tikz}
\usetikzlibrary{calc}

\tikzset{
    right angle quadrant/.code={
        \pgfmathsetmacro\quadranta{{1,1,-1,-1}[#1-1]}     % Arrays for selecting quadrant
        \pgfmathsetmacro\quadrantb{{1,-1,-1,1}[#1-1]}},
    right angle quadrant=1, % Make sure it is set, even if not called explicitly
    right angle length/.code={\def\rightanglelength{#1}},   % Length of symbol
    right angle length=2ex, % Make sure it is set...
    right angle symbol/.style n args={3}{
        insert path={
            let \p0 = ($(#1)!(#3)!(#2)$) in     % Intersection
                let \p1 = ($(\p0)!\quadranta*\rightanglelength!(#3)$), % Point on base line
                \p2 = ($(\p0)!\quadrantb*\rightanglelength!(#2)$) in % Point on perpendicular line
                let \p3 = ($(\p1)+(\p2)-(\p0)$) in  % Corner point of symbol
            (\p1) -- (\p3) -- (\p2)
        }
    }
}
\pagestyle{fancy}
\fancyhf{}
\title{The Reverification of the Reynold's Drag Coefficient of the Free Fall of a Coffee Filter Through the Medium of Air}
\author{Gary Fang}
\date{June 20, 2023}
\rhead{Instructor: Dr. Dykshoorn}
\lhead{Student: Gary Fang}
\chead{Air Resistance Lab Report}
\rfoot{June 20, 2023}
\lfoot{SPH3U7}
\cfoot{Page \thepage \hspace{1pt} of \pageref{LastPage}}
\pagestyle{fancy}

 


\begin{document}

\maketitle
\tableofcontents
\newpage

\section{Overview}
\setcounter{page}{1}
The purpose of this experiment is to find the Reynolds number of air and the drag coefficient of a coffee filter. Detailed instruction can be found on Google Classroom.

\section{Theory}
\label{theory}
\subsection{Terminal Velocity}
For the sake of simplicity, we will assume that the freefall process to be
\begin{enumerate}
    \item Freefall with vertical acceleration of $g$
    \item Reaches terminal velocity, acceleration decreases to 0 immediately
\end{enumerate}
Phase \rom{1}:
\begin{align*}
    d_1&=0\\
    v_1&=0\\
    a&=g
\end{align*}
Phase \rom{2}:
\begin{align*}
    d_{1\text{ (phase \rom{2})}}&=d_{2\text{ (phase \rom{1})}}\\
    v_{1\text{ (phase \rom{2})}}&=v_{2\text{ (phase \rom{1})}}\\
    a&=0
\end{align*}
Now we will derive the terminal velocity $v_d$ in total time $T$ and total height $H$:

\begin{align*}
v_d&=gt_1\label{v=at}\tag{1} \\
H&=\frac 12gt_1^2+v_d\left(T-t_1\right)\label{d=vt+1/2 at^2}\tag{2}
\end{align*}
Isolate $t_1$ from \eqref{v=at}, $t_1=\frac{v_d}{g}$, substitute into \eqref{d=vt+1/2 at^2}:
\begin{align*}
    H&=\frac 12g\left(\frac{v_d}{g}\right)^2+v_d\left(T-\left(\frac{v_d}{g}\right)\right)\\
    0&=\frac{v_d^2}{2g}-v_dT+H\\
    v_d&=\frac{2gT\pm\sqrt{4g^2T^2-4\left(2gH\right)}}{2}\\
    v_d&=gT-\sqrt{g^2T^2-2gH}\label{vd:TH}\tag{3}
\end{align*}
The reason that we rejected the other root is because the maximum possible $v$ (ignore air resistance) is $v=gT$, but the conjugate root adds another radical which is always greater or equal to $0$. $v_d\not>gT$. Therefore we reject the root.

\subsection{Air Resistance}
Assume that the drag for is given by the equation
\begin{align*}
    F_{\text{drag}}&=\alpha v^n=\frac 12C\rho Av^2\label{Fdrag}\tag{4}\\
    mg&=\alpha {v_\text{terminal}}^n\\
    m&=\frac{\alpha}{g}{v_\text{terminal}}^n\\
    \ln m&=n\ln v+\ln{\left(\frac{\alpha}{g}\right)}\label{y=mx+b}\tag{5}
\end{align*}
We can find $n$ and $\alpha$ easily after plotting the above equation.
\begin{align*}
    F_\text{drag}&=\frac 12C\rho Av^n\\
    \alpha &=\frac 12 C\rho A\\
    C&=\frac{2\alpha}{\rho A}=\frac{8ge^B}{\pi\rho d^2}
\end{align*}
For some $B\in\mathbb{R}$, and $d$ is the diameter of the coffee filter



\section{Data Processing}
\subsection{Dimension Analysis}
From \eqref{y=mx+b}, substitute the units:
\begin{align*}
    \ln\left(\text{kg}\right)&=2\ln\left(\text{m/s}\right)+\ln\left(\frac{\alpha \text{s}^2}{\text m}\right)\\
    \ln\left(\text{kg}\right)&=2\ln\left(\text m\right)-2\ln\left(\text s\right)+\ln\left(\alpha\right)+2\ln\left(\text s\right)-\ln\left(\text m\right)\\
    \ln\left(\text{kg}\right)&=\ln\left(\text m\right)+\ln\left(\alpha\right)\\
    \therefore \alpha&=\frac{\text{kg}}{\text m}
\end{align*}
From \eqref{Fdrag}, $\alpha=\frac 12C\rho A$
\begin{align*}
    \frac{\text{kg}}{\text m}&=\frac 12 C\frac{\text{kg}}{\text m^3}\left(\text m^2\right)\\
    \frac{\text{kg}}{\text m}&=C\left(\frac{\text{kg}}{\text m}\right)\\
\end{align*}
$\therefore C$ is unitless


\subsection{Error Propagation of $v_d$}
To minimize the uncertainty, we will take \eqref{vd:TH} and rearrange it:
\begin{align*}
    v_d&=gT-\sqrt{g^2T^2-2gH}\\
    v_d&=gT-\sqrt{g^2T^2-2gH}\left(\frac{gT+\sqrt{g^2T^2-2gH}}{gT+\sqrt{g^2T^2-2gH}}\right)\\
    v_d&=\frac{2gH}{gT+\sqrt{g^2T^2-2gH}}\label{vd}\tag{6}\\
    \delta v_d&=v_d\left(\frac{\delta\left(2gH\right)}{2gH}+\frac{\delta\left(gT+\sqrt{g^2T^2-2gH}\right)}{gT+\sqrt{g^2T^2-2gH}}\right)\\
    \delta v_d&=v_d\left(\frac{2g\delta H}{2gH}+\frac{g\delta T+\delta\left(\sqrt{g^2T^2-2gH}\right)}{gT+\sqrt{g^2T^2-2gH}}\right)\\
    \delta v_d&=v_d\left(\frac{2g\delta H}{2gH}+\frac{g\delta T+\frac 12\left(g^2T^2-2gH\right)^{-\frac 12}\delta\left(g^2T^2-2gH\right)}{gT+\sqrt{g^2T^2-2gH}}\right)\\
    \delta v_d&=\left(\frac{2gH}{gT+\sqrt{g^2T^2-2gH}}\right)\left(\frac{2g\delta H}{2gH}+\frac{g\delta T+\frac 12\left(g^2T^2-2gH\right)^{-\frac 12}\left(2g^2T\delta T+2g\delta H\right)}{gT+\sqrt{g^2T^2-2gH}}\right)\\
\end{align*}
Let $x=gT, y=gH$, so $\delta x=g\delta T, \delta y=g\delta H$, the uncertainty becomes:
\[
    \delta v_d=\left(\frac{2y}{x+\sqrt{x^2-2y}}\right)\left(\frac{\delta y}{y}+\frac{\delta x+\frac 12\left(x^2-2y\right)^{-\frac 12}\left(2x\delta x+2\delta y\right)}{x+\sqrt{x^2-2y}}\right)
\]
After some \st{a ton of} algebraic manipulation, the expression simplifies to 
\begin{align*}
\delta v_d&=\frac{x\delta x+\delta y}{\sqrt{x^2-2y}}-\delta x\\
\delta v_d&=\frac{g^2T\delta T+g\delta H}{\sqrt{g^2T^2-2gH}}-g\delta T\label{vd uncertainty}\tag{7}
\end{align*}
When we graph the \eqref{y=mx+b}, the uncertainty in $x$ will be:
\begin{align*}
    x&=\ln\left(v_d\right)\\
    \delta x&=\frac{\delta v_d}{v_d}\label{delta x}\tag{8}
\end{align*}
Similarly, uncertainty in $y$ will be:
\[\delta y=\frac{\delta m}{m}\label{delta y}\tag{9}\]


\newpage
\subsection{Graphs}
Here are some important constants for this lab:
\begin{enumerate}
    \item $g=9.81$ m/s$^2$
    \item Height $h=(2.15\pm 0.01)$ m
    \item Mass of 20 coffee filters $(17.9\pm 0.1)$ g $\implies$ Mass of a coffee filter $m=\left(8.95\pm0.05\right)\times10^{-4}$ kg
    \item Diameter of a coffee filter $d=(0.150\pm 0.005)$ m
    \item Density of air $\rho=1.225$ kg/m$^3$
\end{enumerate}
\begin{center}
\textit{Table 1. Raw data}
\begin{tabular}{|c|c|c|c|c|}
\hline
    Number of Coffee Filters&	$t_1$ ($\pm$0.05s)&	    $t_2$ ($\pm$0.05s)&	    $t_3$ ($\pm$0.05s)&   Average $t$ ($\pm$0.05s)*\\\hline
    2&   	2.20&	    2.09&	2.37&        	2.22\\\hline
    3&   	2.07&	2.13&	2.03&        	2.08\\\hline
    4&   	1.67&	1.53&	1.80&	        1.67\\\hline
    5&   	1.33&	1.37&	1.58&        	1.43\\\hline
    6&   	1.47&	1.25&	1.48&        	1.40\\\hline
    7&   	1.40&	    1.34&	1.49&        	1.41\\\hline
    8&   	1.40&	    1.28&	1.30&            1.33\\\hline
\end{tabular}
\end{center}
*$t_\text{Average}=\dfrac{t_1+t_2+t_3}{3}$

\begin{center}
    \textit{Table 2. Processed Data}\\
\resizebox{11cm}{!}{
    \begin{tabular}{|c|c|c|c|c|c|}\hline
 $v_d$ \eqref{vd}&	$\delta v_d$ \eqref{vd uncertainty}&	$x$ \eqref{y=mx+b}&$y$ \eqref{y=mx+b}& $\delta x$ \eqref{delta x}&	 $\delta y$ \eqref{delta y}\\\hline
0.9910&	0.0281&	-0.009&	-6.3255&	0.0284&	0.0028\\\hline
1.0630&	0.0321&	0.0611&	-5.9201&	0.0302&	0.0019\\\hline
1.3454&	0.0505&	0.2967&	-5.6324&	0.0375&	0.0014\\\hline
1.5983&	0.0711&	0.4689&	-5.4092&	0.0445&	0.0011\\\hline
1.6328&	0.0743&	0.4903&	-5.2269&	0.0455&	0.0009\\\hline
1.6196&	0.0731&	0.4822&	-5.0728&	0.0451&	0.0008\\\hline
1.7364&	0.0842&	0.5518&	-4.9392&	0.0485&	0.0007\\\hline
    \end{tabular}}
\end{center}

\begin{center}
\begin{tikzpicture}

\begin{axis}[
    enlargelimits=false,
    title={$\ln m$ vs. $\ln v_d$},
    xlabel={$\ln v_d$ No Units},
    ylabel={$\ln m$ No Units},
    xmin=-0.05, xmax=0.65,
    ymin=-7, ymax=-4.8,
    xtick={0,0.1,0.2,0.3,0.4,0.5,0.6},
    ytick={-7,-6,-5,-4},
    legend pos=south west,
    ymajorgrids=true,
    xmajorgrids=true,
    grid style=dashed,
    scale=1.13
]
\addplot[
    color=red,
    only marks,
    mark=square,]
table[meta=y]
{Data.dat};


\addplot [
    domain=-0.1:0.7, 
    samples=2, 
    color=red,
    ]
    {2.1*x-6.21};
\addlegendentry{\(\ln m=n\ln v_d+\ln(\alpha/g)\)}

\addplot [only marks] 
  plot [error bars/.cd, x dir=both, x explicit]
  table [x error plus=delta-x, x error minus=delta-x,
  y error plus = delta-y, y error minus=delta-y
  ]{\mytable};
\end{axis}

\end{tikzpicture}
\end{center}



\subsection{Error Propagation of $n$}
Result from Google Spreadsheet shows the that $n=2.1$, $\ln(\alpha/g)=-6.21$
\begin{align*}
n&=\frac{y_1-y_2}{x_1-x_2}\\
    \frac{\delta n}{n}&=\frac{\delta\left(y_1-y_2\right)}{y_1-y_2}+\frac{\delta\left(x_1-x_2\right)}{x_1-x_2}\\
    \delta n&=n_b\left(\frac{2\delta y_\text{Average}}{y_1-y_2}+\frac{2\delta x_\text{Average}}{x_1-x_2}\right)\\
    \delta n&=2.1\left(\frac{2\sum \delta y}{7(y_1-y_2)}+\frac{2\sum \delta x}{7(x_1-x_2)}\right)\\
    \delta n&=0.3
\end{align*}
Therefore the Reynolds number for air $n=2.1\pm 0.3$ with no units.

\subsection{Error Propagation of $C$}
Let the $y$-intercept $b=-6.21$ From the graph
\begin{align*}
    \ln\left(\frac{\alpha}{g}\right)&=b\\
    \alpha&=ge^b\\
    \delta\alpha&=ge^b\delta b\\
    \delta\alpha&=ge^b\left(\frac{\sum \delta y}{7}\right)\\
    \delta\alpha&=0.00137ge^{-6.21}\\
    C&=\frac{2\alpha}{\rho A}\\
    C&=\frac{8\alpha}{\rho\pi d^2}\\
    C_b&=1.8\\
    \delta C&=C_b\left(\frac{\delta\alpha}{\alpha}+\frac{2\delta d}{d}\right)\\
    \delta C&=C_b\left(\frac{0.00137ge^{-6.21}}{ge^{-6.21}}+\frac{2\cdot 0.005}{0.15}\right)=0.1\\
    \end{align*}
Therefore the drag coefficient of a coffee filter $C=1.8\pm 0.1$ with no units. We would like to express $C$ in the form of
\begin{align*}
    C&=\frac{8ge^B}{\pi\rho d^2}\\
    B&=\ln\left(\frac{\pi\rho d^2C}{8g}\right)=-6.22\\
    \delta B&=\frac{2C\delta d+d\delta C}{d^3C^2}=3.02\\
\end{align*}
Therefore $B=-6.2\pm 3.0$


\pagebreak
\section{Evaluation}
\begin{align*}
    \%\text{Percentage Difference}&=\left|\frac{\text{uncertainty}}{\text{best value}}\right|\times 100\%\\
    \frac{\delta n}{n}&=14\%\\
    \frac{\delta C}{C}&=6\%
\end{align*}
The Reynolds number from the internet is 2
\begin{align*}
    \%\text{Error}&=\left|\frac{\text{expected}-\text{experimental}}{\text{expected}}\right|\times 100\%\\
    E_n&=\frac{2.1-2}{2}\\
    E_n&=5\%
\end{align*}
The drag coefficient of a coffee filter $C=1$
\[E_C=\frac{1.8-1}{1}=80\%\]


\section{Summary}
The experiment shows that the Reynolds number of air is $2.1\pm 0.3$, drag coefficient of a coffee filter is $1.8\pm 0.1$. With percentage difference being 5\% and 80\% respectively. Which is a decent result for $n$, but the result for $C$ requires improvement. To further improve the accuracy of the result, we can record the experiment (dropping the coffee filter) and count the frames instead of using a stopwatch which is unprecise. 

\end{document}